%! AP Statistics — Unit 9, Lesson 9.6 — Group Activity
\documentclass[10pt]{article}
\usepackage{apstats-article}
\usepackage{apstats-boxes}

\begin{document}

\pageheader{Unit 9, Lesson 9.6}{Group Activity}
\namepartnerperiod

\vspace{0.1in}
\textbf{Inference Procedure Identification --- 8 Scenarios}

\textbf{Directions:} Read each scenario. Classify the data and identify the correct inference
procedure at your tier. Work with your group and use the decision framework from notes.

\vspace{0.1in}

\textbf{Scenarios:}
\begin{enumerate}[itemsep=4pt]
  \item A researcher records whether each of 80 randomly selected college students owns a car
        (yes/no). She wants to \textit{estimate} the proportion who own cars.
  \item A school surveys 120 students and records their preferred lunch option (pizza, salad,
        pasta, other). She wants to test whether the four options are equally popular.
  \item A researcher randomly assigns 40 patients to medication A or medication B and records
        recovery time (days) for each. She tests whether the two medications differ in mean
        recovery time.
  \item 50 pairs of identical twins are studied; one twin from each pair exercises 5+ hours per
        week and the other exercises less. The researcher records reaction time for each twin
        and tests whether exercise is associated with faster reaction time.
  \item A political scientist randomly surveys 300 voters in two cities (150 from each) and
        records party affiliation (Democrat, Republican, Independent). She tests whether the
        party affiliation distribution is the same in both cities.
  \item A nutritionist records daily sugar intake (g, $x$) and body fat percentage (\%, $y$)
        for a random sample of 35 adults. She wants to \textit{estimate} the true slope $\beta$.
  \item A researcher surveys 200 people and records education level (HS, Some College,
        Bachelor's, Graduate) and income category (Low, Middle, High) to test whether
        education and income are associated.
  \item A sports scientist randomly samples 60 high school athletes and records their resting
        heart rate (bpm). She wants to test whether the mean resting heart rate differs from
        the national average of 72 bpm.
\end{enumerate}

\vspace{0.1in}

\begin{tcolorbox}[colback=white, colframe=black!40,
    title=\textbf{Tier R --- Remediate},
    fonttitle=\bfseries, arc=2mm, left=3mm, right=3mm, top=2mm, bottom=2mm]

Complete the table for \textbf{Scenarios 1--4} using the decision framework and reference table
from your notes.

\vspace{0.1in}
\renewcommand{\arraystretch}{2.0}
\begin{tabularx}{\linewidth}{@{} c X X X c @{}}
  \rowcolor{sky}
  \textbf{\#} & \textbf{Categorical or Quantitative?} & \textbf{Groups / Structure} & \textbf{Procedure Name} & \textbf{CI or Test?} \\
  \hline
  1 & & & & \\
  2 & & & & \\
  3 & & & & \\
  4 & & & & \\
\end{tabularx}
\end{tcolorbox}

\begin{tcolorbox}[colback=white, colframe=black!40,
    title=\textbf{Tier A --- Approaching Proficiency},
    fonttitle=\bfseries, arc=2mm, left=3mm, right=3mm, top=2mm, bottom=2mm]

Complete the table for \textbf{all 8 scenarios} without a reference table.

\vspace{0.1in}
\renewcommand{\arraystretch}{2.0}
\begin{tabularx}{\linewidth}{@{} c X X X c @{}}
  \rowcolor{sky}
  \textbf{\#} & \textbf{Categorical or Quantitative?} & \textbf{Groups / Structure} & \textbf{Procedure Name} & \textbf{CI or Test?} \\
  \hline
  1 & & & & \\
  2 & & & & \\
  3 & & & & \\
  4 & & & & \\
  5 & & & & \\
  6 & & & & \\
  7 & & & & \\
  8 & & & & \\
\end{tabularx}

\vspace{0.1in}
Then, for scenarios \textbf{1}, \textbf{3}, and \textbf{6}, state the parameter being estimated
or the hypotheses being tested.

\vspace{0.05in}
\textbf{Scenario 1} --- \writeline
\textbf{Scenario 3} --- \writeline
\textbf{Scenario 6} --- \writeline
\end{tcolorbox}

\begin{tcolorbox}[colback=white, colframe=black!40,
    title=\textbf{Tier E --- Extension},
    fonttitle=\bfseries, arc=2mm, left=3mm, right=3mm, top=2mm, bottom=2mm]

Complete all of Tier A above. Then answer the following.

\vspace{0.08in}
\textbf{Bonus 1.} For Scenario 4 (identical twins), explain why a \textit{paired} $t$-test is
more appropriate than a \textit{2-sample} $t$-test.
\writelines{3}

\textbf{Bonus 2.} Scenarios 5 and 7 both involve two categorical variables, yet they use
different chi-square procedures. Explain the key difference between a chi-square test for
\textit{homogeneity} and a chi-square test for \textit{independence}.
\writelines{3}
\end{tcolorbox}

\end{document}
